\documentclass{article}
\usepackage{graphicx} % Required for inserting images
\usepackage{amsmath}
\usepackage{circuitikz}

\title{Op-Amp Problem}
\author{}
\date{}

\begin{document}
\maketitle

We denote the \textit{upper threshold voltage} by $V_h$ and the \textit{lower threshold voltage} by $V_l$. Above $V_h$, the comparator outputs the positive rail voltage $V_s$, and below $V_l$, it outputs the negative rail voltage $-V_s$.\\

\begin{circuitikz}

% Comparator
\draw (6,2) node[op amp, yscale = -1] (opamp) {};

% Output node
\draw (opamp.out) -- (9,2) node[right]{$V_{out}$};

\draw (0.7,2.5) node[left]{$V_{in}$} -- (1,2.5) to[R=$R_1$] (4,2.5);

% --- Hysteresis branch (separate branch from output to + input)
\draw (8,2) -- (8,3.5) to[R=$R_2$] (4,3.5) -- (4,2.5) node[below]{$V_+$} -- (opamp.+);

% --- RC branch (separate branch from output to - input)
\draw (9,2) -- (9,0.5) to[R=$R_3$] (4, 0.5) -- (4,1.5) node[below left]{$V_-$} -- (opamp.-);

% Capacitor to ground from RC node
\draw (4,1) to[C=$C$] (4,-1.7) node[ground]{};

\end{circuitikz}

\section*{Finding $V_h$}
Using voltage division, we can obtain the following expression for $V_+$.

\begin{equation*}
    V_+ = V_{in} + \frac{V_{out} - V_{in}}{R_1+R_2} R_1
\end{equation*}
\noindent
To find an expression for $V_-$ we must solve

\begin{align*}
    I_{R_3} = C \frac{dV_-}{dt} = \frac{V_{out} - V_-}{R_3}\\
    \frac{dV_-}{dt} + \frac{1}{R_3C} V_- = \frac{V_{out}}{R_3C} 
\end{align*}

\noindent
The comparator approaches the positive rail when $V_+ > V_-$, so we substitute both.\\
\hrule

\begin{align*}
    V_{in} - \frac{V_{in} + V_s}{R_f+R} R \ &>\ V_{ref}\\
    \frac{R_f}{R_f+R} V_{in} \ &>\ V_{ref} + \frac{R}{R_f+R} V_s\\
    V_{in} \ &>\ \frac{R_f + R}{R_f} V_{ref} + \frac{R}{R_f} V_s
\end{align*}
\noindent
We have calculated the upper threshold voltage as:
\begin{equation*}
    \boxed{V_h = \frac{R_f + R}{R_f} V_{ref} + \frac{R}{R_f} V_s}
\end{equation*}

\section*{Finding $V_l$}
Now, we assume $V_{out} = V_s$, and again we find an expression for $V_+$.
\begin{equation*}
    V_+ = V_{in} + \frac{V_s - V_{in}}{R_f+R} R
\end{equation*}
\noindent
The comparator approaches the negative rail when $V_+ < V_{ref}$, so we substitute and simplify.
\begin{align*}
    V_{in} + \frac{V_s - V_{in}}{R_f+R} R \ &< \ V_{ref}\\
    \frac{R_f}{R_f+R} V_{in} \ &<\ V_{ref} - \frac{R}{R_f+R} V_s\\
    V_{in} \ &<\ \frac{R_f + R}{R_f} V_{ref} - \frac{R}{R_f} V_s
\end{align*}
\noindent
From that we find the expression for the lower threshold voltage:
\begin{equation*}
    \boxed{V_l = \frac{R_f + R}{R_f} V_{ref} - \frac{R}{R_f} V_s}
\end{equation*}

\end{document}